\paragraph{Automorphic and crystalline realizations at the chosen place.}
Use the coefficient field and the twisted tuple of this
lift at \(w\mid p\). Choose a \(p\)-adic embedding \(j_w\) inducing \(w\) and an isomorphism \(\iota:\overline{\Q}_p\to\C\). On the coefficient field their composite is prescribed:
\[
 \iota\circ j_w=j_{\mathbb C},
\]
where \(j_{\mathbb C}\) is the selected complex embedding. The latter comparison is algebraic.
The finite Hecke operators can be read through it because
their matrices on the finite class model are formed by
the same right translations and transported level
matrices.

At each member of the selected CM half, use
the finite adelic reading extending that \(p\)-adic
embedding. In the adapted basis the level matrices and
their inverses have integral entries. The twist is a
finite character, so adjoining the roots of unity of its
order puts these matrices and character values in the
same ring of integers \(O\). Evaluation at the finite
double classes gives the common integral form model.
The selected tame level makes the stabilizer orders units
in \(O\), including for the smaller weight \(\mu\).
\label{cor:lift-frame-galois-1}

\begin{proposition}[The oriented crystalline frame and the marked string]
\label{prop:lift-frame-crystalline-reading}
In the selected definite model, with the prescribed algebraic
type and spherical local factors over \(p\), the attached
crystalline representation has the rank-\(r\) normalized
Satake readings. For the \(\lambda\) system assume the specified
length-\(t\) divisor of the Hecke polynomial, \(r<2t\), and the
isolated length-\(t\) run in the weight at every label. Its
unique long Speh factor gives the marked consecutive crystalline
string. After the common character twist the \(\mu\)-frame has the following jumps at the two labels, in the given orientation:
\[
 \text{jump at }\tau=h_0-U'_{\tau,\bar a},\qquad\text{jump at }c\tau=U'_{\tau,a}.
\]
The marked string uses this orientation and the Frobenius transport of the proof.
\end{proposition}

\begin{proof}
\proofstep{The Hecke action after transport to the definite group.}
The characteristic-zero coefficient model first gives
a vector-valued form on the original definite space.
The isometry to the standard definite group transports
its rational and level actions. Use the algebraic
coefficient action at infinity to extend the finite
class values. Pair them using the positive invariant
Hermitian form on the coefficient type and take a
nonzero irreducible component of the resulting weight
space, using an invariant orthogonal complement.
The transported finite Hecke sums are the same on these
components.

In the split reading the prime is the
conjugate of the marked row prime. An isometry of the
two split groups identifies the right coset sums using
the dual double-coset label, which sends the list of elementary-divisor exponents to
\[
 (a_1,\ldots,a_r)\longmapsto(-a_r,\ldots,-a_1).
\]
This reads the given
Hecke eigenvalue, under \(\iota\), at the conjugate
prime. At primes of \(E\) above \(p\) the selected
level is spherical and the level character is trivial
on the corresponding hyperspecial group. The operators
indexed by good split primes away from \(p\) include
all coefficients of their Satake polynomials.
\label{cor:lift-frame-galois-2}

\proofstep{Attachment to the normalized Speh factors.}
Use unitary base change
with these local readings, the discrete spectrum
classification for \(\mathrm{GL}_r\), and Galois
attachment for the regular algebraic polarized cuspidal
constituents. Writing \(F=E^+\) for the totally real field of the definite group, the field hypotheses are
\[
 [F:\Q]\ge2,\qquad E/F\text{ unramified at all finite places}.
\] Base change has the
archimedean infinitesimal characters of the selected
algebraic type and the spherical components over \(p\).
Its discrete factors are Speh representations with
conjugate self-dual unitary cuspidal data. The spherical
component of the induced local product forces the
cuspidal local factors to be spherical.

For a cuspidal datum \(\pi\) of rank \(a\) in a Speh factor
of length \(d\), put
\[
 s_j=\frac{r-1}{2}+\frac{d-1}{2}-j,
 \qquad c_j=\frac{a-1}{2}-s_j \qquad (0\le j<d).
\]
In the unitary infinitesimal parameter let \(z\) range over
the exponents of \(\pi\) at a label. Regular algebraicity of
the full rank-\(r\) parameter gives
\[
 s_j-z\in\Z,
\]
because
\(z+(d-1)/2-j\) occurs there. The same regularity makes the
cuspidal exponents distinct at every label. The twist
\(\pi_j=\pi\otimes |\det|^{c_j}\) has exponents \(z+c_j\)
in \((a-1)/2+\Z\), and is polarized by the norm-power
character descending from \(|\cdot|^{2c_j}\). If \(a\) is odd,
multiply the polarizing character on the real subfield by
the quadratic character of the CM extension to give the
required sign at every real place. Its pullback is trivial.

Apply
Barnet-Lamb--Gee--Geraghty--Taylor, Theorem~2.1.1, to this
regular algebraic polarized cuspidal representation
in rank \(a\). With unitary Satake roots \(\alpha_{\pi,h}\) and Hodge labels
\(z_{\pi,\tau,h}\), the reciprocal normalization in rank \(a\)
reads
\[
 \operatorname{HT}_\tau(\rho_{\pi_j})
     =\{(a-1)/2-z_{\pi,\tau,h}-c_j\}_h
     =\{s_j-z_{\pi,\tau,h}\}_h,\qquad
 \operatorname{Sat}(\rho_{\pi_j})=\{q^{s_j}\alpha_{\pi,h}\}_h
\]
at the indicated good split prime for the Satake expression.
The first equality is the Hodge--Tate normalization and the
second multiset uses \((a-1)/2-c_j=s_j\).
The representation is crystalline above \(p\) because
\(\pi_j\) is spherical there. At a good split prime of
norm \(q\), the twist and the rank-\(a\) local Langlands
normalization multiply the unitary Satake roots by
\[
 q^{-c_j+(a-1)/2}=q^{s_j}.
\]

The direct sum over the cuspidal data and \(j\) therefore
has the full rank-\(r\) Satake roots scaled by
\(q^{(r-1)/2}\), with the corresponding full set of
crystalline jumps. Purity in Theorem~2.1.1 gives equal
absolute values for the unramified roots in each
cuspidal factor. If the cuspidal roots in unitary normalization have the common absolute value \(A>0\), then their product has modulus
\[
 A^a=1,
\]
where \(a\) is the rank of this cuspidal factor. This equality uses the unitary central character. Thus \(A=1\).
\cite{labesse-cm,blgg-weight}\label{cor:lift-frame-galois-3}

\proofstep{Identifying the unique long rank-one string.}
For the selected system of weight \(\lambda\) there is
additional information. Its occurrence in the chosen
block of the row supplies, at a good split prime, the
divisor of its Hecke polynomial coming from the
length-\(t\) string. Its center has modulus one in
unitary Satake normalization. To identify the global
factor, each cuspidal Satake multiset in the discrete
decomposition has unit modulus at this prime, by
temperedness of the cohomological conjugate self-dual
cuspidal factors in their unitary normalization. If a Speh constituent has length \(s\) and cuspidal rank \(a\),
its largest root modulus is \(q^{(s-1)/2}\), and its rank
contribution is \(as\). A root of modulus \(q^{(t-1)/2}\)
therefore requires
\[
       s\ge t,\qquad as\le r<2t .
\]
Hence \(a=1\), and two such constituents cannot fit in rank \(r\).

At each
archimedean label the specified length-\(t\) constant
run in the highest weight has strict drops at its two
ends when those ends are present. The shifted
infinitesimal exponents at this label are distinct.
The long Speh constituent gives a consecutive interval
of exponents. It cannot cross either drop, because each
drop leaves an absent intermediate exponent. The other intervals together have at most \(r-s<t\)
entries. They cannot cover the designated consecutive run of
\(t\) entries. The long interval must meet that run and cannot
cross a missing exponent at either boundary, so \(s\le t\).
Combining the size bound for the constituent and the two boundary gaps gives
\[
 t\le s\le t,\qquad s=t.
\]
Thus the long interval is exactly the specified run. Its Galois factors are the \(t\) successive
crystalline twists of a rank-one character. The
remaining Speh factors give \(\rho\).
\label{cor:lift-frame-galois-4}

\proofstep{Filtration jumps for the two CM orientations.}
The common normalization fixes their readings on both
orientations. Write \(\tau\) for a member of the
\(p\)-adic half chosen by the frame, \(c\tau\) for its
conjugate, \(s_\tau\) for the determinant weight shift
there, and
\[
 \bar a=r-1-a
\]
for a reversed coordinate.
Before the common twist, the type on \(\tau\) has highest
weight \(\mu_{\tau,a}+s_\tau\) and the type on \(c\tau\)
has its dual. The algebraic Hecke character with
crystalline jumps \(h_0-r+1+s_\tau\) on \(\tau\) and
\(-s_\tau\) on \(c\tau\) is chosen unramified over \(p\).
The sum of the two jumps is constant on all pairs, as
required for the algebraic character. Under the
reciprocal Frobenius convention the untwisted jump is
\(\bar a-\lambda_{{\rm aut},\bar a}\) in either local
highest weight \(\lambda_{\rm aut}\).

Adding the character
jump gives the two arrays
\[
 J_{\tau,a}=h_0-U'_{\tau,\bar a},\qquad
 J_{c\tau,a}=U'_{\tau,a},\qquad
 U'_{\tau,a}=\mu_{\tau,a}+r-1-a .
\]
Thus the marked orientation for \(U'\) is the conjugate
of the half used for the complex automorphic frame.
Apply the same twist to the selected \(\lambda\) system.
On that orientation the marked line jumps are
\(B_\tau+s\), \(0\le s<t\), with \(B_\tau\) the bottom
of the shifted run of \(\lambda\). On the other orientation
their bottom is \(h_0-B_\tau-(t-1)\), with the order of
the marked indices reversed. Both arrays and both
families of lines therefore lie in \([0,h_0]\).
The Frobenius permutation acts on all embeddings of the
unramified host and commutes with their CM involution.
This determines the transported label and its
orientation together in the crystalline comparison.
\label{cor:lift-frame-galois-5}
\end{proof}
